\chapter{Concepts and Architecture}
\label{ch:architecture}

Blackship is best understood as a transaction coordinator around native
FreeBSD subsystems. It calculates a plan, records ownership, applies host
mutations in order, and rolls them back in reverse order if a later step fails.

\section{Control flow}

\begin{shipfigure}
\begin{tikzpicture}[
  node distance=0.75cm and 0.55cm,
  box/.style={draw=steel,rounded corners=2pt,fill=codeback,minimum height=0.75cm,
              text width=2.3cm,align=center,font=\small},
  arrow/.style={-{Latex[length=2mm]},thick,color=signalred}
]
  \node[box] (cli) {CLI / JSON / Library / MCP};
  \node[box,right=of cli] (manifest) {Manifest and Jailfile parsers};
  \node[box,right=of manifest] (bridge) {Bridge and dependency graph};
  \node[box,below=of manifest] (state) {Scope, audit, network, and lease stores};
  \node[box,below=of cli] (freebsd) {jail, jls, jexec, mount, rc};
  \node[box,below=of bridge] (host) {ZFS, PF, VNET, rctl, cpuset};
  \draw[arrow] (cli) -- (manifest);
  \draw[arrow] (manifest) -- (bridge);
  \draw[arrow] (bridge) -- (host);
  \draw[arrow] (bridge) -- (freebsd);
  \draw[arrow,<->] (bridge) -- (state);
  \draw[arrow] (state) -- (host);
\end{tikzpicture}
\end{shipfigure}

The CLI is intentionally thin. The same core abstractions back ephemeral
library containers and long-lived manifest jails. Errors retain typed causes at
the library boundary and stable codes at the JSON boundary.

\section{Manifest, service, and project names}

The manifest is a desired-state document. Each \setting{[[jails]]} table has a
short service name. Blackship prefixes that name with the project name when it
creates the kernel jail, avoiding collisions between projects. Commands accept
the short name, the full name, an unambiguous prefix, a matching tag, or
\command{ALL} where a selector is allowed.

Dependencies form a directed acyclic graph. Starting a target walks its
prerequisites in topological order. Stopping a target normally walks its
dependents in reverse order. A cycle or unknown dependency is a configuration
error, not an ordering guess.

\section{Releases, roots, and images}

A \emph{release} is a complete FreeBSD userland that can seed a jail. It may
come from official distribution archives, pkgbase, or \command{blackship
commit}. A jail root is a writable instance of that release.

With ZFS, the release dataset owns a \command{@pristine} snapshot and jail
roots are clones. Without ZFS, Blackship copies the release directory. A
Jailfile modifies a root in reproducible steps; on ZFS, each completed
\command{RUN} or \command{COPY} can become a cached layer snapshot.

\section{Lifecycle and state machines}

The lifecycle distinguishes declared intent from partial host reality. A jail
moves through creation, startup, running, draining, stopping, cleanup, and
failure states. Blackship checkpoints ownership before each external mutation.
Consequently, cleanup can reason about a half-created interface or mount even
when no running jail exists.

The main lifecycle operations are:

\begin{description}
  \item[up] provision and start selected jails with their dependencies;
  \item[down] drain and stop selected jails in dependency-safe order;
  \item[restart] perform the stop and start plans;
  \item[eva] reconcile live-safe parameters, resource rules, and CPU placement
    without restarting; and
  \item[gc] compare durable ownership with live host state and reclaim or report
    leftovers.
\end{description}

\subsection*{One chart for lifecycle and health}

Lifecycle and health share a single statechart, so the order in which states
are left and entered is resolved by the chart instead of by hand:

\begin{figure}[htbp]
\centering
\begin{tikzpicture}[
  font=\small,
  leaf/.style={draw=steel,rounded corners=3pt,fill=codeback,minimum width=2.0cm,
               minimum height=0.8cm,align=center},
  group/.style={draw=steel,rounded corners=6pt,inner sep=0.25cm},
  title/.style={font=\footnotesize\bfseries,text=signalred,anchor=north west},
  move/.style={-{Latex[length=2mm]},thick,color=signalred},
  tag/.style={font=\footnotesize\ttfamily,text=voidblack,fill=white,inner sep=1pt}
]
  \node[leaf] (stopped)  at (1.7,5.0) {Stopped};
  \node[leaf] (failed)   at (1.7,1.4) {Failed};
  \node[leaf] (starting) at (5.9,5.0) {Starting};
  \node[leaf] (stopping) at (5.9,1.4) {Stopping};
  \node[leaf] (draining) at (10.6,1.4) {Draining};
  \node[leaf,fill=white] (healthy)   at (9.6,5.6) {Healthy};
  \node[leaf,fill=white] (unhealthy) at (11.7,5.6) {Unhealthy};
  \node[leaf,fill=white] (failing)   at (9.6,4.3) {Failing};
  \node[leaf,fill=white] (suspended) at (11.7,4.3) {Suspended};
  \coordinate (dtop) at (1.7,5.85);
  \coordinate (utop) at (8.0,7.0);
  \coordinate (jtop) at (7.0,7.6);
  \coordinate (rtl) at (8.35,6.7);
  \coordinate (rbr) at (12.95,3.65);

  \begin{scope}[on background layer]
    \node[group,fill=plasma!14,fit=(starting)(stopping)(draining)(utop)(rtl)(rbr)] (up) {};
    \node[group,fill=plasma!14,fit=(stopped)(failed)(dtop)] (down) {};
    \node[group,draw=voidblack,fit=(down)(up)(jtop),inner sep=0.3cm] (jail) {};
    \node[group,fill=ash!55,inner sep=0pt,fit=(rtl)(rbr)] (running) {};
  \end{scope}
  \foreach \bx/\name in {jail/Jail,down/Down,up/Up,running/Running}
    \node[title] at ([shift={(0.15,-0.1)}]\bx.north west) {\name};

  \draw[move] (stopped) -- node[tag,above] {start} (starting);
  \draw[move] (starting.east) -- node[tag,above] {started} (running.west |- starting.east);
  \draw[move] (draining.north |- running.south) -- node[tag,right] {drain} (draining.north);
  \draw[move] ($(running.south west)+(0.4,0)$) -- node[tag,above left] {stop} (stopping.north east);
  \draw[move] (draining.west) -- node[tag,above] {stop} (stopping.east);
  \draw[move] (stopping.north west) -- node[tag,above right,pos=0.45] {stopped} (stopped.south east);
  \draw[move] (failed.north) -- node[tag,right] {recover} (stopped.south);
  \draw[move] (up.west |- failed.east) -- node[tag,above] {fail} (failed.east);
\end{tikzpicture}
\caption{The jail chart. Arrows are the \setting{state-machines} events, with
\texttt{fail} leaving any state inside Up. The four Running leaves move among
themselves with supervisor health verdicts.}
\label{fig:jail-chart}
\end{figure}

Events and guards stay in the \setting{state-machines} jail machine; the chart
only answers which states a move exits (innermost first) and enters (outermost
first). A \command{fail} from a healthy running jail exits Healthy, Running and
Up, then enters Down and Failed. Health verdicts from the supervisor move the
leaf inside Running. When a jail stops, the claims held by Up (the VNET
interface, then the address lease) are released in the order the resolved exit
reports.

\section{Durable stores}

Files beneath the configured data directory serve different purposes:

\begin{shiptable}{L{0.22\textwidth}L{0.68\textwidth}}
\toprule
Store & Role \\
\midrule
Scope & Current ownership and the rollback work-list for a jail. \\
Audit & Append-only JSONL lifecycle history that survives scope deletion. \\
Networks & Persistent named-network definitions, recreated after reboot. \\
Leases & Serialized address ownership for concurrent jail creation. \\
Port forwards & PF mappings that can be reloaded after PF restart. \\
Container registry & Process ownership records for ephemeral library/CLI
containers, including explicit ownership handoff. \\
\bottomrule
\end{shiptable}

Mutable TOML state is staged in a uniquely named temporary file beside its
target and then atomically renamed into place. Keeping the temporary and target
files on one filesystem prevents readers from observing a partially written
scope, network, lease, forward, or container-ownership record.

\section{Failure model}

Host operations can fail between any two syscalls. Blackship therefore avoids
the fiction that a multi-step jail start is atomic. It instead makes rollback
explicit and idempotent. A failed start removes resources it created; a later
\command{cleanup} or \command{gc} can finish the work after interruption.

Resource-limit and VNET attachment failures are fatal. Running a jail after an
isolation control failed would violate the manifest more seriously than
returning an error.

\begin{shipnote}[Inspect before applying]
The structured dry-run plans are calculated from the same dependency and drift
logic used by real operations. Use \command{--format json ... --dry-run} in a UI
or \command{--dry-run} at the terminal before broad changes.
\end{shipnote}
